% La cuenta del contador de bits tras cada una de las cuatro llamadas del ejemplo.
% Clave ficticia de 40 hexadecimales. Posiciones calculadas con engine.bits.BitsEngine.
\begin{tikzpicture}[x=0.345cm,y=0.5cm,font=\ttfamily\scriptsize]
  \def\clave{3,f,9,a,1,c,0,7,e,2,b,4,d,8,5,a,6,c,0,f,1,7,e,9,3,a,b,2,d,4,c,8,8,1,5,e,7,f,6,0}
  % #1 fila (y), #2 candidatas, #3 confirmadas, #4 rótulo izquierdo, #5 rótulo derecho
  \newcommand{\fila}[5]{%
    \foreach \i in {0,...,39} {\draw[ceniza!45,line width=0.25pt] (\i,#1) rectangle ++(1,1);}
    \foreach \i in {#2} {\fill[oroclaro!55] (\i,#1) rectangle ++(1,1);
                         \draw[oro,line width=0.4pt] (\i,#1) rectangle ++(1,1);}
    \foreach \i in {#3} {\ifnum\i>-1 \fill[tinta] (\i,#1) rectangle ++(1,1);\fi}% -1 = ninguna
    \foreach \c [count=\i from 0] in \clave {%
      \gdef\tiracolor{ceniza!75}%
      \foreach \j in {#2} {\ifnum\i=\j \gdef\tiracolor{tinta}\fi}%
      \foreach \j in {#3} {\ifnum\i=\j \gdef\tiracolor{white}\fi}%
      \node[inner sep=0pt,text=\tiracolor] at (\i+0.5,#1+0.5) {\c};}
    \node[anchor=east,font=\sffamily\footnotesize,text=tinta] at (-0.5,#1+0.5) {#4};
    \node[anchor=west,font=\sffamily\footnotesize,text=tinta] at (40.5,#1+0.5) {#5};
  }
  \fila{0}{0,1,2,3}{-1}{1}{16 bits}
  \fila{-1.6}{0,...,7}{-1}{2}{32 bits}
  \fila{-3.2}{0,...,7,12,13,14,15}{-1}{3}{48 bits}
  \fila{-4.8}{0,...,7,12,13,14,15}{16,...,27}{4}{96 bits}
  \foreach \i in {0,10,20,30,39} {\node[font=\sffamily\tiny,text=ceniza] at (\i+0.5,1.45) {\i};}
  % leyenda
  \fill[oroclaro!55] (0,-6.6) rectangle ++(1,1); \draw[oro,line width=0.4pt] (0,-6.6) rectangle ++(1,1);
  \node[anchor=west,font=\sffamily\footnotesize,text=tinta] at (1.4,-6.1) {candidata (tramo de 4 a 7 caracteres)};
  \fill[tinta] (20,-6.6) rectangle ++(1,1);
  \node[anchor=west,font=\sffamily\footnotesize,text=tinta] at (21.4,-6.1) {confirmada (tramo de 8 o más)};
\end{tikzpicture}
